% Part: turing-machines
% Chapter: undecidability
% Section: introduction

\documentclass[../../../include/open-logic-section]{subfiles}

\begin{document}

\olfileid[fr]{tur}{und}{int}
\olsection{Introduction}

Il peut sembler évident que toutes les fonctions ne sont pas
calculables, même parmi les fonctions arithmétiques : elles
sont trop nombreuses, et certaines ont un comportement trop
complexe. Songeons, par exemple, à des fonctions définies
à partir de la désintégration de particules radioactives,
ou d'autres phénomènes chaotiques ou aléatoires. Fixons un
instant initial et comptons les intervalles successifs d'une
seconde. Définissons $f(n)$ comme le nombre de particules
de l'univers qui se désintègrent pendant le $n$-ième de
ces intervalles. Voilà, semble-t-il, une fonction dont
on ne saurait espérer calculer les valeurs.

Ne pas imaginer comment calculer une telle fonction est
cependant bien différent de démontrer qu'elle n'est pas
calculable. Même une fonction qui paraît d'une complexité
insurmontable peut être calculable au sens abstrait.
Supposons, par exemple, que l'univers ait une durée finie
et finisse, dans un avenir extrêmement lointain, par se
contracter en un point. Il n'y aurait alors qu'un nombre
fini, quoique prodigieux, d'intervalles d'une seconde après
l'instant initial pour lesquels $f(n)$ serait définie.
Or toute fonction définie sur un ensemble fini d'entrées
est calculable : en principe, ses valeurs pourraient figurer
dans une immense table ou être codées dans le diagramme
de transition d'une très grande machine de Turing.
Cette existence abstraite n'exige pas que nous connaissions
effectivement la table.

Nous nous intéressons souvent aux fonctions dont les
valeurs répondent à des questions par oui ou par non.
Par exemple, à la question « $n$ est-il un nombre premier ? »
correspond la fonction
\[
\fn{isprime}(n) = \begin{cases}
  1 & \text{si $n$ est premier}\\
  0 & \text{sinon.}
  \end{cases}
\]
On dit qu'une question de ce type est
\emph{effectivement décidable} si la fonction associée,
à valeurs dans $\{0,1\}$, est effectivement calculable.

Pour démontrer mathématiquement que certaines fonctions
ne sont pas effectivement calculables ou que certains
problèmes ne sont pas effectivement décidables, il est
indispensable de fixer un modèle précis de calcul, puis
d'établir que ce modèle ne calcule pas ces fonctions ou
ne décide pas ces problèmes. Nous pouvons ainsi montrer
que les machines de Turing ne calculent pas toutes les
fonctions et ne décident pas tous les problèmes. La thèse
de Church--Turing permet ensuite de conclure qu'aucune
procédure effective ne peut combler cette limite.

L'argument décisif repose sur la possibilité d'attribuer
des nombres aux machines de Turing elles-mêmes. Le moyen
le plus simple est de les énumérer : fixons une façon
d'écrire les machines et leurs programmes, puis dressons-en
une liste systématique. Dès que cette énumération est
établie, un argument de cardinalité montre qu'il existe
des fonctions non calculables par machine de Turing.
En effet, l'ensemble des fonctions de $\Nat$ vers~$\Nat$
--- et même celui des fonctions de $\Nat$ vers
$\{0,1\}$ --- est !!{nonenumerable}, tandis que l'ensemble
des machines de Turing, donc celui des fonctions qu'elles
calculent, est seulement !!{denumerable}.

Nous pouvons aussi définir des fonctions et des problèmes
\emph{particuliers} dont nous démontrons respectivement
la non-calculabilité et l'indécidabilité. Le
\emph{problème de l'arrêt} en est un exemple. Une machine
de Turing possède une description finie : la liste de
ses instructions. Cette description, autrement dit son
programme, peut elle-même servir d'entrée à une autre
machine de Turing. On peut donc poser à une machine
des questions sur d'autres machines. L'une d'elles
est : « La machine décrite finit-elle par s'arrêter
sur l'entrée~$n$ ? » Une machine capable de répondre
serait bien pratique, comme un dispositif de contrôle
détectant les boucles infinies. Mais Turing a démontré
qu'aucune machine de Turing ne peut décider cette question.
Il n'existe pas de machine unique qui, recevant la
description d'une machine~$M$ et un nombre~$n$,
s'arrête toujours avec la sortie~$1$ ou~$0$ selon
que~$M$ s'arrête ou non sur l'entrée~$n$.

Une fois certains problèmes particuliers reconnus
indécidables, nous pouvons nous en servir pour en
établir d'autres. Un exemple célèbre est la question :
« La !!{formula} du premier ordre~$!A$ est-elle valide ? »
Aucune machine de Turing ne peut, pour toute
!!{formula} du premier ordre~$!A$ reçue en entrée,
s'arrêter à coup sûr avec la sortie~$1$ ou~$0$
selon que~$!A$ est valide ou non. Historiquement,
la recherche d'une procédure effective pour répondre
à cette question s'appelait simplement le
« problème de la décision »; nous disons donc que
ce problème est insoluble. Turing et Church ont
démontré ce résultat indépendamment, à peu près
à la même époque. On l'appelle aussi le
théorème de Church--Turing.\footnote{Cette appellation, conservée de la source, désigne ici le résultat d'indécidabilité de la validité du premier ordre. Il faut la distinguer de la thèse de Church--Turing, qui relie une notion intuitive de calcul effectif à un modèle formel.}

\end{document}
